\section{Recursive functions and the fixed machine model}
\label{sec:annex-p04-recursive-machine-equivalence}

\begin{annexitem}{annex:p04.recursive-counter-realization}{thm:p04-recursive-counter-realization}{Counter constructions}
\begin{proof}
We use only increment, zero-test/decrement, and finite control, as in
\cref{def:p04-counter-program}. The zeroing of a counter is a loop which
decrements it until its zero branch is taken. An unconditional transfer
is the zero branch of a fresh counter kept at zero. Each description
below therefore expands into a finite list of the permitted instructions.

First specify a copying block. For distinct source and destination
counters, zero the destination and a fresh temporary counter. Until the
source is zero, decrement it and increment both the destination and the
temporary. If its original value was \(a\), after \(j\) iterations their
values are \(a-j,j,j\). Then decrement the temporary while incrementing
the source. This restores the source to \(a\), leaves the destination at
\(a\), and clears the temporary. Both loops terminate. Fresh scratch
counters keep all other named counters unchanged.

Addition to a counter is repeated increment controlled by a copied
argument; multiplication is repeated addition controlled by another
copy. Comparison copies both arguments and decrements the copies in
lockstep until one is zero. Testing the first copy before decrementing
distinguishes \(a\leq b\) from \(a>b\), including equality at zero.
Subtraction of a smaller counter value is a bounded sequence of
decrements. These descriptions terminate and preserve any inputs whose
copies they use; they appeal to no machine--recursion equivalence.

The fixed pairing and its inverse can also be handled by these blocks.
To encode \((a,b)\), form \(u=a+b\), accumulate
\(1+2+\cdots+u\), and add \(b\). To decode \(z\), start with
\(u=t=0\), and while \(t+u+1\leq z\), replace
\((u,t)\) by \((u+1,t+u+1)\). At the loop boundary
\(t=u(u+1)/2\leq z\). It terminates, since the next triangular
number eventually exceeds \(z\), and gives
\[
 b=z-t,\qquad a=u-b.
\]
The stopping condition implies \(0\leq b\leq u\), so these are natural
numbers and their pairing is \(z\). Iterating the inverse a fixed
number of times decodes the fixed-arity tuple code. This supplies the
input counters from the procedure's single numerical input counter.

For the closure induction it is useful to give every block designated
input counters, a distinct result counter, and one exit. On return it
preserves its input counters and places its value in the result counter.
Scratch counters are disjoint from those used by surrounding blocks.
There are only finitely many of them for any fixed construction.
A block's entry initializes the scratch values it needs; hence repeated
entries do not rely on their values from an earlier call.

The zero block clears its result counter. The successor block copies its
input and increments the copy. A projection copies the selected input.
These have the specified behavior and always return.

For composition, run the blocks for all inner functions, one after
another, on the unchanged input counters, putting their results in
distinct counters. Only after every inner block has returned, run the
outer block on those results. If any required inner block fails to
return, execution remains in it; if the outer block fails to return,
execution remains there. Otherwise the returned value is the prescribed
composition. Thus the construction is strict even when the outer
function would ignore an argument whose inner computation diverges.

For primitive recursion in the manuscript's last-argument convention,
\[
 f(\vec x,0)=g(\vec x),\qquad
 f(\vec x,n+1)=h(\vec x,n,f(\vec x,n)),
\]
keep the input counters unchanged and copy \(n\) into a countdown
counter. Initialize an index counter to zero, and run the block for
\(g\) to obtain the current value. If the countdown is zero, exit.
Otherwise decrement it, copy the current value into a separate argument
counter, and run the block for \(h\) on the parameters, the index, and
that value. Replace the current value by the returned result, increment
the index, and repeat. At the boundary after \(i\) completed steps, the
index is \(i\), the countdown is \(n-i\), and the current value is
\(f(\vec x,i)\). Induction establishes this invariant whenever that
boundary is reached. A divergence in the base or a required step blocks
every subsequent boundary, exactly as the definition requires. If all
required values exist, there are exactly \(n\) step calls and the block
returns the required value. Inputs remain unchanged.

For minimization, keep \(\vec x\) unchanged and set a candidate counter
to zero. Run the block for \(g(\vec x,y)\). If it returns zero, copy
\(y\) to the result and exit. If it returns a nonzero value, increment
\(y\) and repeat. Reaching candidate \(y\) means that every earlier
value was defined and nonzero. Thus a return has exactly the least-zero
meaning in \cref{def:p04-partial-recursive}. If an earlier required
computation diverges, execution stays in that call and never skips it.
If all tested values are defined and nonzero, there are infinitely many
iterations and no return. These exhaust the undefined cases.

The least class closed under these finitary operations is the union of
its finite construction stages: that union contains the initial
functions and is itself closed under each operation. Every function
under discussion therefore has a finite construction tree. Instantiate
the displayed finite blocks along that tree, rename scratch counters and
instruction locations to avoid clashes, prepend fixed-arity tuple
decoding, and replace the final exit by return of the result counter.
The result is a finite counter procedure. The argument proves its
behavior by induction on the construction tree, independently of any
decision about eventual termination of its execution.
\end{proof}
\end{annexitem}

\begin{annexitem}{annex:p04.fixed-machine-pr}{thm:p04-fixed-machine-pr}{Finite simulation with a fixed origin}
\begin{proof}
We first give the elementary primitive-recursive closure facts used in
the coding, rather than infer them from Turing computability.
Predecessor, addition, multiplication, and truncated subtraction
\(u\mathbin{\dot{-}}v=\max(u-v,0)\) are obtained by primitive recursion.
Zero-test is obtained by a recursion with initial value one and all
successor values zero. These give sign, equality, order, and conditional
selection. For a zero--one test \(e\), selection between \(u\) and \(v\)
is \(eu+(1\mathbin{\dot{-}}e)v\).
Bounded sums and products are primitive recursions. A bounded search
for the least \(j\leq n\) satisfying a primitive-recursive test iterates
through \(0,\ldots,n\), retaining its first successful index and using
\(n+1\) as the failure value. The retained index is the state of a
primitive recursion; once it is at most \(n\), subsequent updates leave
it unchanged. Thus the search is primitive recursive.
If a scalar recursion needs a nullary initial constant, first carry an
unused positive-arity parameter and then identify or discard the dummy
input by composition. All functions used below can consequently be
formed in the manuscript's positive-arity partial-recursive convention
as well as in its primitive-recursive class.

For \(d>0\), the quotient is explicitly
\[
 \left\lfloor n/d\right\rfloor
   =\sum_{j=1}^{n}\mathbf1_{jd\leq n};
\]
give it value zero at \(d=0\). Remainder is obtained by multiplication
and truncated subtraction. Powers are primitive recursive by recursion
on the exponent. Triangular numbers are the recurrence
\(t(0)=0,\ t(u+1)=t(u)+u+1\), so the fixed pairing
\(t(a+b)+b\) is primitive recursive without an initial division
operation. Its inverse finds the largest \(u\leq z\) with
\(t(u)\leq z\), and returns \(z-t(u)\) and
\(u-(z-t(u))\) in the appropriate order. The largest such \(u\) is
\(\sum_{j=1}^{z}\mathbf1_{t(j)\leq z}\). Since
\(t(u)\leq z<t(u+1)\), the two recovered components are nonnegative
and pair back to \(z\). Fixed tuple coding and decoding are therefore
primitive recursive by finite composition. No unbounded search has
entered these closure arguments.

Fix the state and symbol numbers of
\cref{def:p04-origin-preserving-code}. Every nonnegative integer \(W\)
has a unique base-\(b\) digit sequence with finite nonzero support, so
it describes a finite-support tape under the address bijection. Write
\[
 D(W,j)=\left\lfloor W/b^j\right\rfloor\bmod b.
\]
This is primitive recursive. The address-code operations for a move
right or left are respectively
\[
 r(j)=\begin{cases}j+2,&j\text{ even},\\
                   j\mathbin{\dot{-}}2,&j\text{ odd},\end{cases}
 \qquad
 \ell(j)=\begin{cases}1,&j=0,\\
             j\mathbin{\dot{-}}2,&j>0\text{ even},\\
             j+2,&j\text{ odd}.\end{cases}
\]
Checking the cases \(i\geq0\) and \(i<0\), including the crossings
at zero and minus one, proves
\(r(\operatorname{addr}(i))=\operatorname{addr}(i+1)\) and
\(\ell(\operatorname{addr}(i))=\operatorname{addr}(i-1)\).
Parity is a remainder test, so both maps are primitive recursive.

Decode \(c\) into \((q,j,W)\). If \(q\) is a valid nonhalting state,
read \(a=D(W,j)\), and use the fixed finite transition table to obtain
\((q',a',d)\). Finite case distinction on the state and symbol numbers
is primitive recursive. Replace the tape code by
\[
 W'=\bigl(W\mathbin{\dot{-}}a b^j\bigr)+a'b^j,
\]
and replace \(j\) by \(r(j)\) or \(\ell(j)\) according to \(d\).
Here \(W\geq ab^j\); digit expansion shows that this removes exactly
the scanned digit and inserts exactly the new one. All other digits
are unchanged. Pair the new components to obtain \(N(c)\).
If the state is halting or is an invalid state number, define \(N(c)=c\).
This is a total primitive recursive function, and the preceding
calculation proves its exact transition meaning on configuration codes.

For initialization let \(x=\operatorname{code}_k(\vec x)\).
The length of \(\operatorname{bin}(x)\) is one if \(x=0\), and
otherwise the least positive \(v\leq x+1\) with \(x<2^v\).
This bounded search always succeeds. If the length is \(v\), its
\(i\)-th bit, most significant first, is
\[
 e_i=\left\lfloor x/2^{v-1-i}\right\rfloor\bmod2
       \quad(0\leq i<v).
\]
These formulas also give the single zero bit when \(x=0\).
The initial tape code is
\[
 W_0=\sum_{i<v}(1+e_i)b^{2i}.
\]
The digit at address \(\operatorname{addr}(i)=2i\) is therefore exactly
the symbol number of the input bit, and all other cells are blank.
Pair \((q_0,0,W_0)\) to obtain \(I(\vec x)\).
The length, bits, bounded sum, and final pairing are primitive recursive.

For output, find the first \(m\leq W+1\) with \(D(W,2m)=0\).
There is such an \(m\): for \(e>W\), the elementary inequality
\(b^e\geq e+1>W\) gives \(D(W,e)=0\); in particular the cell with
address code \(2(W+1)\) is blank. Thus bounded search finds exactly
the first blank at a nonnegative physical address. The prefix is valid
precisely when \(m>0\), every \(D(W,2i)\) with \(i<m\) is one or two,
and its first digit is two unless \(m=1\). These tests express both
binary-symbol validity and the prohibition of leading zero in a
multi-symbol output. They are primitive recursive by bounded products
and the elementary tests above. Let \(G(c)\) be their conjunction
with the test that the state is the accepting state. In the valid case
the output value is
\[
 O(c)=\sum_{i<m}
       \bigl(D(W(c),2i)\mathbin{\dot{-}}1\bigr)2^{m-1-i}.
\]
Use that same bounded expression as a total default when the prefix is
invalid. It is primitive recursive and has the claimed value whenever
\(G(c)=1\). The test does not inspect the current head position, and
digits outside the prefix do not affect the output. This is exactly the
fixed numerical interface, not the different head-relative delimiter
convention used elsewhere for a separately specified coding.

Finally, primitive recursion on the supplied step count gives
\(C(\vec x,n)\). Induction on \(n\) shows that it codes the actual
computation up to a halt and thereafter repeats the halting code.
Every operation in this construction was obtained by the displayed
primitive-recursive closures; none uses the theorem being proved.
\end{proof}
\end{annexitem}

\begin{annexitem}{annex:p04.tm-partial-recursive}{thm:p04-tm-partial-recursive}{The two directions}
\begin{proof}
If \(f\) is partial recursive, \cref{thm:p04-recursive-counter-realization}
gives a finite counter procedure receiving its fixed tuple code and
returning exactly its defined values. Apply
\cref{thm:p04-counter-compilation}. Its explicit one-tape construction
decodes the most-significant-bit-first input, simulates the counters,
and writes the canonical binary output at the preserved physical origin.
Every finite counter instruction takes finitely many machine steps;
a nonreturning counter run becomes an infinite machine run. Hence the
compiled machine meets both clauses of \cref{def:partial-computable}.
The counter-compilation proof uses finite transition templates and
their invariants, not the present equivalence theorem.

Conversely suppose a machine computes \(f\) in that definition's exact
sense. Take the primitive recursive \(C,G,O\) from
\cref{thm:p04-fixed-machine-pr}, and form
\[
 t(\vec x)=\mu n\,[1\mathbin{\dot{-}}G(C(\vec x,n))=0].
\]
Every tested value is defined, being primitive recursive. Therefore
\(t\) is defined exactly when a valid accepting configuration is
reached, and then it is its first time. By the witness machine's
specification this occurs exactly when \(f(\vec x)\) is defined;
on other inputs its computation is infinite. Consequently
\[
 f(\vec x)\simeq O(C(\vec x,t(\vec x))),
\]
where equality includes simultaneous undefinedness. The right side is
partial recursive by minimization and strict composition. This proves
the other direction under the same tuple and output encodings.
No identification with an index from a separately chosen normal-form
coding is needed.
\end{proof}
\end{annexitem}

\begin{annexitem}{annex:p04.primitive-recursive-total}{thm:p04-primitive-recursive-total}{The total subclass}
\begin{proof}
Zero, successor, and projections are total. Composition of total
functions is total, and primitive recursion on a natural-number argument
makes finitely many calls to its total base and step functions. Induction
on a primitive-recursive construction proves totality. The counter
construction above uses only these terminating base, composition, and
recursion blocks when no minimization occurs, so its compiled machine
halts with the required value on every input. A nullary subconstruction
is a fixed natural constant and can be returned by literal initialization;
equivalently it can be carried with an unused positive-arity parameter.
For the nullary function itself, the required input is
\(\operatorname{bin}(\operatorname{code}_0())=\operatorname{bin}(0)\);
a machine writing the fixed constant in canonical binary at cell zero
computes it. Thus the zero-arity convention introduces no exception.
\end{proof}
\end{annexitem}
